%%%%%%%%%%%%%%%%%%%%%%%%%%%%%%%%%%%%%%%%%%%%%%%%%%%%%%%%%%%%%%%%%%%%%%
%
% Example of \TeX Format for AQIS Pre-Proceedings
%
%%%%%%%%%%%%%%%%%%%%%%%%%%%%%%%%%%%%%%%%%%%%%%%%%%%%%%%%%%%%%%%%%%%%%%%

\documentclass[a4paper]{article}
\usepackage{aqis}

\begin{document}

\title{
   Example of \TeX\ Format for AQIS Pre-Proceedings
}

\author{
  AuthorA
  \affiliation{1}
   \affiliation{2}
  \email{username1@domainname1}
  \and
 AuthorB
   \affiliation{2}
  \email{username2@domainname2}
  \and
 AuthorC
   \affiliation{3}
  \email{username3@domainname3}
  }
  
\address{1}{
  Please input affiliation 1 here
}
\address{2}{
  Please input affiliation 2 here
}
\address{3}{
  Please input affiliation 3 here
}


\abstract{
Please substitute this for your 100-words abstracts.
}

\keywords{AQIS, template}


\maketitle

%########### body

\section{Format of the Manuscripts}
The camera-ready version should be submitted
via the camera-ready submission page.
It should follow the style described below.

\paragraph{Basic Format:}

\begin{itemize}
  \item
      When you create PDF file from your manuscript source
      (especially from MS-Word), please select fonts carefully for
      some fonts cannot be handled well by PDF creating system.
  \item
    The manuscript should
    not exceed three pages including the title page,
    using A4-sized papers.
  \item
    No page of the manuscript should contain page numbers/page headings.
\end{itemize}

\paragraph{Layout:}

\begin{itemize}
  \item
    The manuscript should
    have top-, foot-, and side-margins of 2cm, 2.5cm, and 1.5cm,
    respectively (the title page should have top margin of 3cm).
  \item
    The manuscript should begin with a title,
    followed by names, affiliations and postal addresses of authors,
    followed by a 100-word abstract and keywords.
    The main body should use two-column style.
    E-mail addresses of authors should be placed in footnotes.
\end{itemize}

\section{Theorems and Proofs}

Our style file prepares the following environments.

\begin{center}
\begin{tabular}{ll}
\hline
Environment name & Heading\\
\hline
{\tt definition} & {\bf Definition}\\
{\tt theorem} & {\bf Theorem}\\
{\tt lemma} & {\bf Lemma}\\
{\tt corollary} & {\bf Corollary}\\
{\tt proposition} & {\bf Proposition}\\
{\tt example} & {\bf Example}\\
{\tt remark} & {\bf Remark}\\
{\tt fact} & {\bf Fact}\\
{\tt conjecture} & {\bf Conjecture}\\
\hline
\end{tabular}
\end{center}

\section{Figures and Tables}

For figures and tables, the captions should be placed
below figures and above tables.

\section{Bibliography Styles}

The following bibliography style used in this document is preferable.
The authors may use common abbreviations for journal and conference names.

\begin{thebibliography}{9}

\bibitem{Grover96}
L.~K.~Grover.
\newblock A fast quantum mechanical algorithm for database search.
\newblock In {\em Proc.\ of the 28th ACM STOC}, pages 212--219, 1996.

\bibitem{Gruska99}
J.~Gruska.
\newblock {\em Quantum Computing}.
\newblock McGraw-Hill, 1999.

\bibitem{Holevo96}
A.~S.~Holevo.
\newblock The capacity of quantum channel with general signal states.
\newblock quant-ph/9611023, 1996.

\bibitem{Sho97SIComp}
P.~W.~Shor.
\newblock Polynomial-time algorithms for prime factorization
and discrete logarithms on a quantum computer.
\newblock {\em SIAM J. on Comp.}, 26(5):1484--1509, 1997.

\end{thebibliography}

\end{document}
% end of file
